\documentclass[14pt,a4paper]{extarticle}
\usepackage{fontspec}
\setmainfont{times.ttf}[Path=C:/Windows/Fonts/,BoldFont=timesbd.ttf,ItalicFont=timesi.ttf,BoldItalicFont=timesbi.ttf]
\setmonofont{consola.ttf}[Path=C:/Windows/Fonts/,BoldFont=consolab.ttf,ItalicFont=consolai.ttf,BoldItalicFont=consolaz.ttf]
\newfontfamily\cyrillicfonttt{cour.ttf}[Path=C:/Windows/Fonts/,BoldFont=courbd.ttf,ItalicFont=couri.ttf,BoldItalicFont=courbi.ttf]
\usepackage{polyglossia}
\setdefaultlanguage{russian}
\setotherlanguage{english}
\usepackage[left=30mm,right=15mm,top=20mm,bottom=20mm]{geometry}
\usepackage{setspace}
\onehalfspacing
\usepackage{indentfirst}
\setlength{\parindent}{1.25cm}
\setlength{\parskip}{0pt}
\usepackage{amsmath,amssymb,array,tabularx,longtable}
\usepackage{float,graphicx,tikz,listings,caption,tocloft,titlesec,fancyhdr,hyperref}
\usetikzlibrary{arrows.meta,shapes.geometric,positioning,calc}
\hypersetup{unicode=true,colorlinks=false,pdfborder={0 0 0},pdftitle={Лабораторная работа № 3. Разворот числа},pdfauthor={Понарин А. А.}}
\urlstyle{same}
\captionsetup{labelsep=endash,font=normalsize}
\captionsetup[figure]{name=Рисунок,justification=centering}
\captionsetup[table]{name=Таблица,justification=raggedright,singlelinecheck=false,margin=0pt}
\captionsetup[lstlisting]{justification=raggedright,singlelinecheck=false,margin=5pt}
\renewcommand{\lstlistingname}{Листинг}
\renewcommand{\contentsname}{СОДЕРЖАНИЕ}
\renewcommand{\cftsecleader}{\cftdotfill{\cftdotsep}}
\renewcommand{\cfttoctitlefont}{\hfill\bfseries\large}
\renewcommand{\cftaftertoctitle}{\hfill}
\setlength{\cftbeforesecskip}{3pt}
\titleformat{\section}{\bfseries\large}{\thesection}{0.5em}{}
\titleformat{\subsection}{\bfseries}{\thesubsection}{0.5em}{}
\titleformat{\subsubsection}{\bfseries}{\thesubsubsection}{0.5em}{}
\titlespacing*{\section}{1.25cm}{0pt}{1em}
\titlespacing*{\subsection}{1.25cm}{1em}{0.6em}
\titlespacing*{\subsubsection}{1.25cm}{1em}{0.6em}
\pagestyle{fancy}\fancyhf{}\fancyfoot[C]{\thepage}
\renewcommand{\headrulewidth}{0pt}\setlength{\headheight}{17pt}
\lstdefinestyle{cpp}{language=C++,basicstyle=\ttfamily\footnotesize,keywordstyle=\mdseries,columns=fullflexible,keepspaces=true,showstringspaces=false,breaklines=true,frame=single,framerule=0.3pt,xleftmargin=3mm,xrightmargin=3mm,aboveskip=0.6em,belowskip=0.6em}
\lstset{style=cpp}
\newcommand{\structhead}[1]{\clearpage\phantomsection\addcontentsline{toc}{section}{#1}\begin{center}\bfseries #1\end{center}\vspace{0.5em}}
\newcommand{\conclusion}[1]{\par\medskip\noindent\begin{minipage}{\linewidth}\setlength{\parindent}{1.25cm}\noindent\hspace*{\parindent}\textbf{Вывод}\par\indent #1\par\end{minipage}\par}
\newcommand{\fig}[3]{\begin{figure}[H]\centering\input{figures/#1.tex}\caption{#2}\label{#3}\end{figure}}
\numberwithin{equation}{section}
\numberwithin{figure}{section}
\numberwithin{table}{section}
\AtBeginDocument{\counterwithin{lstlisting}{section}}
\setlength{\emergencystretch}{2em}
\clubpenalty=10000
\widowpenalty=10000
\displaywidowpenalty=10000
\begin{document}
\begin{titlepage}\thispagestyle{empty}
\noindent\begin{minipage}[c]{0.16\textwidth}\includegraphics[width=\linewidth]{b_logo.jpg}\end{minipage}\hfill
\begin{minipage}[c]{0.82\textwidth}\centering\small\bfseries
Министерство науки и высшего образования Российской Федерации\\
Федеральное государственное автономное образовательное учреждение\\
высшего образования\\
«Московский государственный технический университет имени Н.~Э.~Баумана\\
(национальный исследовательский университет)»\\
(МГТУ им. Н.~Э.~Баумана)
\end{minipage}
\par\vspace{0.5em}\noindent\rule{\textwidth}{2pt}\par\vspace{1.5em}
\noindent ФАКУЛЬТЕТ \underline{«Информатика и системы управления»}\\[1em]
{\small\noindent КАФЕДРА \underline{«Программная инженерия»}}
\vspace{2.2cm}
\begin{center}\bfseries\Large Лабораторная работа № 3\\[0.3em]
по дисциплине «Анализ алгоритмов»\end{center}
\vspace{2.2cm}
\noindent\textbf{Тема} \underline{Рекурсивные и нерекурсивные алгоритмы}\\[1em]
\noindent\textbf{Студент} \underline{Понарин А.~А.}\\[1em]
\noindent\textbf{Группа} \underline{ИУ7-55Б}\\[1em]
\noindent\textbf{Преподаватели} \underline{Волкова Л.~Л., Строганов Д.~В.}
\vfill\begin{center}2026~г.\end{center}
\end{titlepage}
\clearpage\pagenumbering{arabic}\setcounter{page}{2}
\renewcommand{\contentsname}{СОДЕРЖАНИЕ}
\tableofcontents

\structhead{ВВЕДЕНИЕ}
Цель работы --- провести сравнение рекурсивного и нерекурсивного алгоритмов на основе графовых моделей алгоритмов.

Для достижения поставленной цели необходимо решить следующие задачи:
\begin{enumerate}
\item дать определения четырёх графовых моделей алгоритмов, операционного и информационного отношений;
\item разработать нерекурсивный и рекурсивный последовательные алгоритмы разворота числа согласно варианту;
\item реализовать разработанные алгоритмы;
\item провести тестирование реализаций алгоритмов;
\item описать каждый алгоритм четырьмя графовыми моделями;
\item обосновать, какие фрагменты алгоритмов или их реализаций могут выполняться параллельно, а какие должны выполняться последовательно.
\end{enumerate}

\clearpage\section{Аналитическая часть}
\subsection{Постановка задачи}
Дано натуральное число $n$, десятичная запись которого не содержит нулей. Требуется получить целое число, записанное теми же цифрами в противоположном порядке. Например, при $n=179$ результат равен $971$. Функция должна возвращать число целиком; вывод отдельных цифр вместо формирования результата не допускается.

В рекурсивном решении нельзя использовать циклы, строки, списки, массивы и срезы. Допустимы рекурсивные вызовы и целочисленная арифметика. Для выполнения общей задачи сравнения дополнительно разрабатывается нерекурсивный алгоритм с циклом. Таким образом, ограничение варианта соблюдается рекурсивным решением, а итеративное решение служит требуемым объектом сравнения.

Пусть
\begin{equation}
n=\sum_{i=0}^{k-1} c_i10^i,\qquad c_i\in\{1,\ldots,9\}.
\end{equation}
Искомая функция имеет вид
\begin{equation}
\operatorname{rev}(n)=\sum_{i=0}^{k-1}c_i10^{k-1-i}.
\end{equation}
В реализации принимаются числа от $1$ до $10^{19}-1$ без нулевых цифр. Это ограничение машинного представления: все такие числа и их развороты помещаются в 64-разрядный беззнаковый целый тип. Отрицательные числа, ноль и числа с нулевыми цифрами не принадлежат области определения.

\subsection{Операционное и информационное отношения}
Рассматривается последовательное выполнение программы~\cite{voevodin}. \textit{Срабатывание оператора} --- одно конкретное исполнение оператора; один оператор цикла или рекурсивной функции может срабатывать многократно.

\textit{Операционное отношение} связывает два срабатывания, если второе исполняется непосредственно после первого в выбранном последовательном выполнении. Оно описывает порядок исполнения, включая вычисление условий и передачу управления.

\textit{Информационное отношение} связывает срабатывание, выработавшее значение, со срабатыванием, которое использует именно это значение как операнд. Между записью и чтением не должно быть другой записи, заменившей передаваемое значение. Это прямая зависимость по данным, а не всякая последовательность действий с одинаковым именем переменной.

При переходе от срабатываний к операторам связь считается возможной, если она возникает хотя бы в одном допустимом выполнении. Для рекурсии различаются передача аргументов и получение результата: результат вызова становится доступен после возврата из соответствующего дочернего вызова.

\subsection{Четыре графовые модели}
Используются две компактные модели и две истории конкретного выполнения~\cite{voevodin,parallel-course}. Графы ориентированные; стрелка направлена от предшествующего действия к последующему либо от производителя значения к потребителю.

\textit{Граф управления} содержит вершину для каждого оператора, включая условия. Дуга обозначает возможное непосредственное следование операторов. Повторения цикла и рекурсии отображаются возвратными дугами.

\textit{Информационный граф} содержит вершины операторов программы, включая проверки условий. Дуга обозначает возможную прямую информационную связь между их срабатываниями. Проверка условия получает значения переменных как операнды, поэтому к ней проводятся входящие информационные дуги. Выбор дальнейшего пути исполнения отражается в графе управления. У компактного графа возможны циклы и петли: они могут связывать разные исполнения одного оператора.

\textit{Операционная история} строится для заданного входа. Каждому срабатыванию оператора соответствует отдельная вершина, а дуги задают непосредственный последовательный порядок срабатываний. Для завершившегося однопоточного выполнения она является конечной направленной цепью.

\textit{Информационная история} содержит отдельные вершины срабатываний операторов и прямые информационные связи, возникшие при данном выполнении. Она не содержит ориентированных циклов: значение передаётся от уже выполненного действия последующему. В отличие от компактного информационного графа, история различает итерации и уровни рекурсии.

Вершины передачи параметров и возврата результата далее включаются как граничные операторы передачи данных. Переход к компактной модели выполняется отождествлением вершин историй, соответствующих одному месту исходного кода.

\subsection{Условия параллельного выполнения}
Отсутствие пути по данным между двумя срабатываниями позволяет рассматривать их совместное исполнение при выполненных условиях управления. При работе с общей памятью дополнительно проверяются конфликты чтения и записи: нельзя уничтожить ещё не прочитанное значение. Например, два чтения одного числа не конфликтуют.

Информационная история показывает доступный параллелизм вычислений, но сама по себе не разрешает удалять проверки условий. Перестановка действий должна сохранять допустимость арифметики и результат программы. Анализ в четвёртой части относится к исходным последовательным алгоритмам, а не к измеренной многопоточной реализации.
\conclusion{Для сравнения способов разворота числа используются четыре графовые модели. Граф управления и операционная история отражают порядок выполнения, а информационный граф и информационная история --- зависимости по данным, определяющие возможности параллельного выполнения.}

\clearpage\section{Конструкторская часть}
\subsection{Нерекурсивный алгоритм}
Используется накопитель $a$, первоначально равный нулю. На каждой итерации выделяются последняя цифра $d=n\bmod10$ и оставшаяся часть $q=\lfloor n/10\rfloor$. Цифра присоединяется справа к накопителю: $a\gets10a+d$. Затем $n\gets q$. При $n=0$ возвращается накопитель.

Схема нерекурсивного алгоритма приведена на рисунке~\ref{fig:flow-i}.
\fig{flow_iterative}{Схема нерекурсивного алгоритма разворота числа}{fig:flow-i}

\clearpage\subsection{Рекурсивный алгоритм}
Вспомогательная функция $R(n,a)$ принимает ещё не обработанную часть числа и накопленный результат. Она определяется соотношением
\begin{equation}
R(n,a)=\begin{cases}
a,&n=0,\\
R(\lfloor n/10\rfloor,10a+n\bmod10),&n>0.
\end{cases}
\end{equation}
Первый вызов --- $R(n,0)$. Базовый случай соответствует $n=0$. При $n>0$ выполняется следующий рекурсивный вызов. Его результат сохраняется в локальной переменной $t$ и возвращается вызывающей функции. Все параметры передаются по значению; разделяемого изменяемого накопителя нет.

Схема рекурсивного алгоритма приведена на рисунке~\ref{fig:flow-r}.
\fig{flow_recursive}{Схема рекурсивного алгоритма разворота числа}{fig:flow-r}

\clearpage\subsection{Обоснование корректности и завершения}
Пусть исходное число имеет цифры $c_{k-1}\ldots c_1c_0$. После обработки $j$ цифр, $0\le j\le k$, выполняется инвариант
\begin{equation}
n_j=\left\lfloor\frac{n_0}{10^j}\right\rfloor,
\qquad
a_j=\sum_{i=0}^{j-1}c_i10^{j-1-i},
\end{equation}
где при $j=0$ сумма равна нулю. В исходном состоянии $n_0$ не изменено, $a_0=0$, поэтому инвариант верен.

При переходе к следующему шагу вычисляется $d=c_j$. Умножение $a_j$ на 10 и прибавление $c_j$ дают запись $c_0\ldots c_j$, а целочисленное деление $n_j$ на 10 удаляет обработанную цифру. Следовательно, инвариант сохраняется.

После $k$ шагов $n_k=0$, а $a_k=\operatorname{rev}(n_0)$. При нерекурсивном вычислении возвращается значение $a_k$. При рекурсивном вычислении это значение возвращается из базового случая и затем без изменения передаётся через все ожидающие вызовы.

На каждом шаге при $n>0$ выполняется $0\le\lfloor n/10\rfloor<n$, поэтому число шагов конечно. Для $k$ цифр выполняются $k$ итераций либо $k$ рекурсивных вызовов при $n>0$ и один вызов при $n=0$. В реализации глубина вычислительной рекурсии не превышает 20 вызовов.

Переполнение исключено: $a_j\le10^j-1$. При $j<k\le19$ очередные $10a_j$ и $10a_j+c_j$ не превышают $10^{19}-1<2^{64}-1$.

\subsection{Требования к программе}
Программа предоставляет режимы выбора одного или обоих алгоритмов со сравнением результатов. Вводится одно число в строке; пробелы по краям допустимы. Знаки, нулевые цифры, посторонние символы и несколько чисел отклоняются.

Функции возвращают целое число без ввода-вывода и вещественной арифметики. Аргумент проверяется до вычисления: рекурсивно для рекурсивной версии, итеративно для нерекурсивной.
\conclusion{Оба способа разворота корректны и завершаются; в выбранном диапазоне переполнение исключено.}

\clearpage\section{Технологическая часть}
\subsection{Средства реализации}
Для реализации выбран язык C++17 со статической типизацией~\cite{stroustrup}. Для представления чисел используется тип \texttt{std::uint64\_t}.

Применён компилятор GCC версии 14.2.0 (MSYS2 UCRT64). Параметры компиляции:
\begin{itemize}\setlength{\itemsep}{0pt}\setlength{\parsep}{0pt}
\item \texttt{-std=c++17} --- выбор версии языка C++17;
\item \texttt{-Wall -Wextra -Wpedantic} --- включение предупреждений компилятора;
\item \texttt{-O2} --- оптимизация программы.
\end{itemize}

\subsection{Программная реализация алгоритмов}
\noindent\begin{minipage}{\linewidth}
\setlength{\parindent}{1.25cm}
В листинге~\ref{lst:i} приведён исходный код функции нерекурсивного разворота числа.

\lstinputlisting[firstline=14,lastline=23,caption={Функция нерекурсивного разворота числа},label={lst:i}]{algorithms.cpp}
\end{minipage}

\medskip\noindent\begin{minipage}{\linewidth}
\setlength{\parindent}{1.25cm}
В листинге~\ref{lst:r} приведён исходный код функции рекурсивного разворота числа. При первом вызове параметр \texttt{a} равен нулю. Циклы, строки и массивы не используются.

\lstinputlisting[firstline=25,lastline=34,caption={Функция рекурсивного разворота числа},label={lst:r}]{algorithms.cpp}
\end{minipage}

\clearpage\subsection{Функциональное тестирование}
Тестовые данные включают следующие категории:
\begin{itemize}
\item положительные --- числа с различными цифрами;
\item граничные --- минимальное и максимальное допустимые числа;
\item специальные --- одинаковые цифры и палиндром;
\item отрицательные --- ноль, нулевые цифры, неверный формат и превышение допустимого диапазона.
\end{itemize}
Результаты приведены в таблице~\ref{tab:tests}. Обозначение «Ошибка» соответствует отклонению ввода с диагностическим сообщением.

\begin{table}[H]\centering
\caption{Результаты функционального тестирования}\label{tab:tests}
\small\setlength{\tabcolsep}{3pt}
\begin{tabularx}{\textwidth}{|*{3}{>{\centering\arraybackslash}X|}}\hline
\bfseries Входные данные & \bfseries Ожидаемый результат & \bfseries Фактический результат\\\hline\hline
1 & 1 & 1\\\hline
12 & 21 & 21\\\hline
179 & 971 & 971\\\hline
11111 & 11111 & 11111\\\hline
12321 & 12321 & 12321\\\hline
1234567891234567891 & 1987654321987654321 & 1987654321987654321\\\hline
9999999999999999999 & 9999999999999999999 & 9999999999999999999\\\hline
0 & Ошибка & Ошибка\\\hline
10 & Ошибка & Ошибка\\\hline
101 & Ошибка & Ошибка\\\hline
$-179$ & Ошибка & Ошибка\\\hline
1.5 & Ошибка & Ошибка\\\hline
abc & Ошибка & Ошибка\\\hline
12 34 & Ошибка & Ошибка\\\hline
Пустая строка & Ошибка & Ошибка\\\hline
11111111111111111111 & Ошибка & Ошибка\\\hline
18446744073709551615 & Ошибка & Ошибка\\\hline
99999999999999999999 & Ошибка & Ошибка\\\hline
\end{tabularx}\end{table}

\conclusion{При функциональном тестировании обеих реализаций все тестовые данные обработаны с ожидаемым результатом.}

\clearpage\section{Исследовательская часть}
Графовые модели соответствуют функциям разворота числа из листингов~\ref{lst:i} и~\ref{lst:r}. Проверки ввода, диагностика и выбор режима в них не рассматриваются. Каждый оператор присваивания считается одной вершиной; умножение и сложение внутри $10a+d$ не дробятся на машинные команды.

Вершины обозначены номерами операторов. В историях один номер повторяется при каждом срабатывании соответствующего оператора. Показано семейство историй для произвольного допустимого числа $n$ из $k$ цифр. Многоточия обозначают пропущенные повторения; при малом $k$ лишние повторения отсутствуют.

\subsection{Графовые модели нерекурсивного алгоритма}
Соответствие вершин исходному коду из листинга~\ref{lst:i} приведено в таблице~\ref{tab:inodes}.
\begin{table}[H]\centering\caption{Вершины моделей нерекурсивного алгоритма}\label{tab:inodes}
\small\begin{tabularx}{\textwidth}{|c|>{\raggedright\arraybackslash}X|}\hline
\bfseries Вершина & \bfseries Оператор / действие\\\hline
0 & Получение исходного $n$\\\hline
1 & $a\gets0$\\\hline
2 & Проверка $n\ne0$\\\hline
3 & $d\gets n\bmod10$\\\hline
4 & $q\gets\lfloor n/10\rfloor$\\\hline
5 & $a\gets10a+d$\\\hline
6 & $n\gets q$\\\hline
7 & Возврат $a$\\\hline
\end{tabularx}\end{table}

Граф управления и информационный граф приведены на рисунках~\ref{fig:ic} и~\ref{fig:ii} соответственно. Проверка 2 получает значение $n$ из вершины 0 либо 6.
\fig{iter_control}{Граф управления нерекурсивного алгоритма}{fig:ic}
\fig{iter_information}{Информационный граф нерекурсивного алгоритма}{fig:ii}

Операционная и информационная истории для числа из $k$ цифр приведены на рисунках~\ref{fig:io} и~\ref{fig:id}. После $k$ обработанных цифр проверка 2 даёт ложный результат и выполняется возврат 7.
\fig{iter_operational}{Операционная история нерекурсивного алгоритма для $k$ цифр}{fig:io}
\fig{iter_dynamic}{Информационная история нерекурсивного алгоритма для $k$ цифр}{fig:id}

\subsection{Графовые модели рекурсивного алгоритма}
Соответствие вершин исходному коду из листинга~\ref{lst:r} приведено в таблице~\ref{tab:rnodes}. Передача аргументов и получение результата дочернего вызова обозначены разными вершинами 6 и 7.
\begin{table}[H]\centering\caption{Вершины моделей рекурсивного алгоритма}\label{tab:rnodes}
\small\begin{tabularx}{\textwidth}{|c|>{\raggedright\arraybackslash}X|}\hline
\bfseries Вершина & \bfseries Оператор / действие\\\hline
0 & Получение параметров $(n,a)$ текущего вызова\\\hline
1 & Проверка $n=0$\\\hline
2 & Возврат $a$ в базовом случае\\\hline
3 & $d\gets n\bmod10$\\\hline
4 & $q\gets\lfloor n/10\rfloor$\\\hline
5 & $b\gets10a+d$\\\hline
6 & Передача аргументов $(q,b)$ дочернему вызову\\\hline
7 & Получение результата дочернего вызова в $t$\\\hline
8 & Возврат $t$ из текущего вызова\\\hline
\end{tabularx}\end{table}

Граф управления и информационный граф приведены на рисунках~\ref{fig:rc} и~\ref{fig:ri}. Проверка 1 получает $n$ из параметров в вершине 0. Возврат выполняется в продолжение соответствующего родительского вызова; локальные переменные разных глубин различаются.
\fig{rec_control}{Граф управления рекурсивного алгоритма}{fig:rc}
\fig{rec_information}{Информационный граф рекурсивного алгоритма}{fig:ri}

Операционная и информационная истории для числа из $k$ цифр приведены на рисунках~\ref{fig:ro} и~\ref{fig:rd}. После $k$ вызовов при $n>0$ выполняется базовый вызов при $n=0$. Перед получением результата в вершине 7 завершается дочерний вызов; один и тот же результат передаётся обратно через все уровни.
\fig{rec_operational}{Операционная история рекурсивного алгоритма для $k$ цифр}{fig:ro}
\fig{rec_dynamic}{Информационная история рекурсивного алгоритма для $k$ цифр}{fig:rd}

\subsection{Анализ возможности параллельного выполнения}
Независимые действия и ограничения для функций из листингов~\ref{lst:i} и~\ref{lst:r} приведены в таблице~\ref{tab:parallel}. Номера вершин соответствуют таблицам~\ref{tab:inodes} и~\ref{tab:rnodes}.
\begin{table}[H]\centering\caption{Независимые действия и зависимости}\label{tab:parallel}
\small\begin{tabularx}{\textwidth}{|p{4.8cm}|>{\raggedright\arraybackslash}X|}\hline
Фрагмент & Возможность параллельного выполнения\\\hline
Обе функции: 3 и 4 & Вычисления $d$ и $q$ независимы: общий вход $n$, разные выходы\\\hline
Нерекурсивная: 5 и 6 & Обновления $a$ и $n$ независимы после получения $d$, $q$ и чтения старого $n$\\\hline
Рекурсивная: 4 и 5 & Вычисления $q$ и $b$ независимы при готовом $d$\\\hline
Вершина 5 соседних итераций & Последовательны: следующий накопитель зависит от предыдущего\\\hline
Рекурсивная: 6 и 0 дочернего вызова & Дочерний вызов начинается после вычисления его аргументов $q$ и $b$\\\hline
Возвраты соседних глубин & Последовательны: родителю требуется результат потомка\\\hline
\end{tabularx}\end{table}

\subsection{Сравнение графовых моделей и ресурсов}
Сопоставление реализаций приведено в таблице~\ref{tab:compare}.
\begin{table}[H]\centering\caption{Сопоставление реализаций}\label{tab:compare}
\small\begin{tabularx}{\textwidth}{|>{\raggedright\arraybackslash}X|c|c|}\hline
Свойство & Нерекурсивная & Рекурсивная\\\hline
Организация повторений & Цикл & Вызовы функции\\\hline
Время выполнения при $k$ цифрах & $\Theta(k)$ & $\Theta(k)$\\\hline
Дополнительная память & $\Theta(1)$ & $O(k)$\\\hline
Вершины операционной истории & $5k+4$ & $8k+3$\\\hline
Вершины информационной истории & $5k+4$ & $8k+3$\\\hline
Прямые информационные дуги & $6k+2$ & $10k+2$\\\hline
Критический путь & $\Theta(k)$ & $\Theta(k)$\\\hline
Циклы и массивы & Один цикл & Отсутствуют\\\hline
\end{tabularx}\end{table}

Здесь $k$ --- число цифр. Критический путь определяется цепочкой из $k$ обновлений накопителя. Память рекурсии оценена для обычного стека вызовов; стоимость операции с машинным целым считается постоянной.
\conclusion{Построены восемь графовых моделей. Установлено, что частное и остаток, а также отдельные присваивания могут вычисляться независимо. Цепочка накопителя и передача результатов рекурсии ограничивают параллелизм; длина критического пути у обоих алгоритмов линейна по числу цифр.}

\structhead{ЗАКЛЮЧЕНИЕ}
Цель работы достигнута: проведено сравнение рекурсивного и нерекурсивного алгоритмов разворота числа на основе графовых моделей алгоритмов.

В ходе работы решены следующие задачи:
\begin{itemize}
\item даны определения четырёх графовых моделей алгоритмов, операционного и информационного отношений;
\item разработаны нерекурсивный и рекурсивный последовательные алгоритмы разворота числа согласно варианту;
\item реализованы разработанные алгоритмы;
\item проведено функциональное тестирование обеих реализаций;
\item каждый алгоритм описан графом управления, информационным графом, операционной и информационной историями;
\item определены фрагменты исходного кода, допускающие параллельное выполнение, и последовательные зависимости.
\end{itemize}

При обоих способах разворота числа трудоёмкость линейна по числу цифр. Для нерекурсивной реализации требуется постоянная дополнительная память, для рекурсивной с обычным стеком вызовов --- линейная. Частное и остаток от деления на 10 могут вычисляться параллельно. Обновления накопителя и передача результатов между уровнями рекурсии сохраняют последовательные зависимости; длина критического пути у обеих реализаций составляет $\Theta(k)$, где $k$ --- число цифр.

\clearpage\phantomsection
\addcontentsline{toc}{section}{СПИСОК ИСПОЛЬЗОВАННЫХ ИСТОЧНИКОВ}
\begingroup
\renewcommand{\refname}{СПИСОК ИСПОЛЬЗОВАННЫХ ИСТОЧНИКОВ}
\titleformat{name=\section,numberless}{\centering\bfseries}{}{0pt}{}
\titlespacing*{\section}{0pt}{0pt}{1em}
\begin{thebibliography}{9}\sloppy
\bibitem{voevodin} Воеводин, В.~В., Воеводин, В.~В. Параллельные вычисления [Текст] / В.~В.~Воеводин, В.~В.~Воеводин. --- СПб.: БХВ, 2002. --- 608~с.
\bibitem{parallel-course} Параллельная обработка данных: программа курса // Информационно-аналитический центр по параллельным вычислениям PARALLEL.RU. --- URL: \url{https://parallel.ru/info/parallel/parallel_programme.html} (дата обращения: 26.10.2026).
\bibitem{stroustrup} Stroustrup B. The C++ Programming Language. --- 4th ed. --- Addison-Wesley, 2013. --- ISBN 978-0-321-56384-2.
\end{thebibliography}
\endgroup
\end{document}
